\section{CƠ SỞ LÝ THUYẾT}

\subsection{Cloud Computing}

\subsubsection{Tổng quan Cloud Computing}
Cloud Computing (Điện toán đám mây) là mô hình cung cấp tài nguyên máy tính theo yêu cầu qua mạng Internet, cho phép truy cập thuận tiện, theo yêu cầu đến một nhóm tài nguyên máy tính có thể cấu hình được (mạng, máy chủ, lưu trữ, ứng dụng và dịch vụ) mà có thể được cung cấp và giải phóng nhanh chóng với nỗ lực quản lý tối thiểu hoặc tương tác với nhà cung cấp dịch vụ.

Theo Viện Tiêu chuẩn và Công nghệ Quốc gia Hoa Kỳ (NIST), Cloud Computing có 5 đặc điểm thiết yếu:

\begin{itemize}
    \item \textbf{On-demand self-service (Tự phục vụ theo yêu cầu):} Người dùng có thể tự động cấp phát khả năng tính toán như thời gian máy chủ và lưu trữ mạng khi cần thiết mà không cần tương tác con người với từng nhà cung cấp dịch vụ.
    
    \item \textbf{Broad network access (Truy cập mạng rộng):} Các khả năng có sẵn trên mạng và được truy cập thông qua các cơ chế tiêu chuẩn thúc đẩy việc sử dụng bởi các nền tảng khách hàng khác nhau (điện thoại di động, máy tính bảng, laptop, workstation).
    
    \item \textbf{Resource pooling (Tổng hợp tài nguyên):} Tài nguyên máy tính của nhà cung cấp được tổng hợp để phục vụ nhiều người dùng sử dụng mô hình multi-tenant, với các tài nguyên vật lý và ảo khác nhau được gán động và gán lại theo nhu cầu của người dùng.
    
    \item \textbf{Rapid elasticity (Khả năng co giãn nhanh):} Các khả năng có thể được cung cấp và giải phóng một cách linh hoạt, trong một số trường hợp tự động, để mở rộng nhanh chóng ra ngoài và thu hẹp nhanh chóng vào trong tương xứng với nhu cầu.
    
    \item \textbf{Measured service (Dịch vụ đo lường):} Hệ thống cloud tự động kiểm soát và tối ưu hóa việc sử dụng tài nguyên bằng cách tận dụng khả năng đo lường ở một số mức độ trừu tượng phù hợp với loại dịch vụ (lưu trữ, xử lý, băng thông).
\end{itemize}

\subsubsection{Các mô hình dịch vụ}
Cloud Computing cung cấp ba mô hình dịch vụ chính:

\paragraph{IaaS (Infrastructure as a Service - Hạ tầng như Dịch vụ)}
IaaS cung cấp các tài nguyên tính toán cơ bản như máy chủ ảo, lưu trữ, mạng và hệ điều hành. Người dùng có toàn quyền kiểm soát hệ điều hành, ứng dụng và cấu hình mạng, nhưng không quản lý hoặc kiểm soát cơ sở hạ tầng cloud bên dưới.

\textbf{Đặc điểm:}
\begin{itemize}
    \item Cung cấp: Virtual Machines, Storage, Network, Load Balancers
    \item Người dùng quản lý: OS, Middleware, Runtime, Applications
    \item Nhà cung cấp quản lý: Hardware, Virtualization, Networking
    \item Ví dụ: Amazon EC2, Microsoft Azure Virtual Machines, Google Compute Engine, DigitalOcean
\end{itemize}

\paragraph{PaaS (Platform as a Service - Nền tảng như Dịch vụ)}
PaaS cung cấp một nền tảng cho phép khách hàng phát triển, chạy và quản lý các ứng dụng mà không cần phải xây dựng và duy trì cơ sở hạ tầng thường cần thiết cho việc phát triển và khởi chạy ứng dụng.

\textbf{Đặc điểm:}
\begin{itemize}
    \item Cung cấp: Runtime environment, Database, Web Server, Development Tools
    \item Người dùng quản lý: Applications, Data
    \item Nhà cung cấp quản lý: OS, Middleware, Runtime, Infrastructure
    \item Ví dụ: Google App Engine, Heroku, Microsoft Azure App Service, AWS Elastic Beanstalk
\end{itemize}

\paragraph{SaaS (Software as a Service - Phần mềm như Dịch vụ)}
SaaS cung cấp các ứng dụng hoàn chỉnh chạy trên cơ sở hạ tầng cloud. Người dùng truy cập ứng dụng từ nhiều thiết bị khách hàng thông qua giao diện thin client như trình duyệt web.

\textbf{Đặc điểm:}
\begin{itemize}
    \item Cung cấp: Application hoàn chỉnh
    \item Người dùng quản lý: Chỉ cấu hình ứng dụng
    \item Nhà cung cấp quản lý: Toàn bộ stack công nghệ
    \item Ví dụ: Gmail, Microsoft Office 365, Salesforce, Dropbox, Zoom
\end{itemize}

\subsubsection{Các mô hình triển khai}
Cloud Computing có bốn mô hình triển khai:

\begin{enumerate}
    \item \textbf{Public Cloud:} Cơ sở hạ tầng cloud được cung cấp cho việc sử dụng mở cho công chúng. Nó có thể được sở hữu, quản lý và vận hành bởi một tổ chức kinh doanh, học thuật, hoặc chính phủ, hoặc một sự kết hợp của chúng. Ví dụ: AWS, Microsoft Azure, Google Cloud Platform.
    
    \item \textbf{Private Cloud:} Cơ sở hạ tầng cloud được cung cấp để sử dụng độc quyền bởi một tổ chức duy nhất bao gồm nhiều người dùng. Nó có thể được sở hữu, quản lý và vận hành bởi tổ chức, bên thứ ba, hoặc một sự kết hợp của chúng.
    
    \item \textbf{Hybrid Cloud:} Cơ sở hạ tầng cloud là sự kết hợp của hai hoặc nhiều cơ sở hạ tầng cloud riêng biệt (private, public) vẫn là các thực thể duy nhất, nhưng được ràng buộc với nhau bằng công nghệ tiêu chuẩn hoặc độc quyền cho phép tính di động của dữ liệu và ứng dụng.
    
    \item \textbf{Community Cloud:} Cơ sở hạ tầng cloud được cung cấp để sử dụng độc quyền bởi một cộng đồng cụ thể của người dùng từ các tổ chức có chung mối quan tâm.
\end{enumerate}

\subsection{Software Defined Networking (SDN)}
\subsubsection{Tổng quan về SDN}

SDN là mô hình mạng dựa trên phần mềm, trong đó tách biệt mặt điều khiển (Control Plane) khỏi mặt chuyển tiếp dữ liệu (Data Plane).

Trong SDN, SDN controller là nơi tập trung các chức năng điều khiển, các thiết bị chuyển mạch chuyển tiếp gói tin theo các luật do Controller thiết lập.

Nhờ SDN, toàn bộ mạng được giám sát và quản lý tập trung thông qua phần mềm và linh hoạt điều chỉnh theo nhu cầu và chính sách của doanh nghiệp.

\subsubsection{Kiến trúc của SDN}

Kiến trúc SDN được chia thành ba lớp chính : Lớp ứng dụng (Application Layer), lớp điều khiển (Control Layer) và lớp cơ sỏ hạ tầng (Infrastructure Layer), liên kết với nhau thông qua giao diện lập trình (API)
\begin{enumerate}
  \item Application Layer
  
  Application Layer là lớp cao nhất của SDN, bao gồm các ứng dụng điều khiển mạng như : routing, firewall, QoS,...
  
  Các ứng dụng gửi yêu cầu đến SDN Controller thông qua API để Controller tính toán và cài đặt các flow rule phù hợp
  
  \item Control Layer 
  
  Cotrol Layer là lớp trung tâm của SDN, bao gồm các SDN Controller : ONOS, OpenDaylight, Ryu, Flloodlight.

  Controller thực hiện các nhiệm vụ như : thu thập thông tin toàn mạng, tính toán đường đi tối ưu, giám sát thiết bị, thiết lập các Flow rule,...và duy trì cái nhìn tổng thể, cho phép đưa ra các quyết định tối ưu cho hệ thống.
  
  \item Infrastructure Layer
  
  Infrastructure Layer gồm các thiết bị thực hiện chuyển tiếp dữ liệu dựa trên các Flow rule đã được Controller cài đặt.
\end{enumerate}
\subsubsection{Cơ chế hoạt động của SDN}
\begin{enumerate}
    \item \textbf{Gói tin đến Switch}
    \begin{itemize}
        \item Switch nhận gói tin từ thiết bị đầu cuối.
        \item Kiểm tra \textbf{Flow Table} để tìm quy tắc xử lý.
    \end{itemize}

    \item \textbf{Kiểm tra Flow Table}
    \begin{itemize}
        \item \textbf{Có Flow Rule:} Switch xử lý ngay theo quy tắc (chuyển tiếp, loại bỏ, sửa đổi gói tin,...).
        \item \textbf{Không có Flow Rule:} Switch gửi thông điệp \textbf{Packet-In} đến SDN Controller.
    \end{itemize}

    \item \textbf{Controller xử lý}
    \begin{itemize}
        \item Phân tích thông tin gói tin (MAC, IP, giao thức, cổng,...).
        \item Xem xét các chính sách mạng (QoS, bảo mật, tải mạng,...).
        \item Tính toán đường đi hoặc hành động phù hợp.
    \end{itemize}

    \item \textbf{Cài đặt Flow Rule}
    \begin{itemize}
        \item Controller gửi thông điệp \textbf{Flow-Mod} đến Switch.
        \item Switch cập nhật \textbf{Flow Table} với quy tắc mới, bao gồm:
        \begin{itemize}
            \item \textbf{Match Fields:} Điều kiện so khớp (MAC, IP, VLAN, cổng, giao thức,...).
            \item \textbf{Actions:} Hành động (chuyển tiếp, loại bỏ, sửa đổi, gửi lên Controller,...).
            \item \textbf{Priority:} Mức độ ưu tiên.
            \item \textbf{Timeout:} Thời gian hiệu lực.
        \end{itemize}
    \end{itemize}

    \item \textbf{Chuyển tiếp gói tin}
    \begin{itemize}
        \item Switch thực hiện xử lý theo Flow Rule mới.
        \item Các gói tin tiếp theo của cùng một \textit{flow} được xử lý trực tiếp tại Switch mà không cần gửi lại Controller, giúp giảm độ trễ và tăng hiệu năng.
    \end{itemize}
\end{enumerate}

\subsection{VLAN}

VLAN (Virtual Local Area Network - Mạng cục bộ ảo) là công nghệ cho phép phân chia một mạng LAN vật lý thành nhiều mạng LAN logic độc lập. Mỗi VLAN tạo ra một broadcast domain riêng biệt, cho phép các thiết bị trong cùng VLAN giao tiếp với nhau như thể chúng được kết nối vào cùng một switch vật lý, bất kể vị trí vật lý thực tế của chúng.

\textbf{Lợi ích của VLAN:}
\begin{itemize}
    \item \textbf{Bảo mật:} Phân đoạn mạng theo phòng ban, giới hạn truy cập giữa các nhóm
    \item \textbf{Giảm broadcast traffic:} Mỗi VLAN là một broadcast domain riêng, giảm tải cho mạng
    \item \textbf{Quản lý linh hoạt:} Dễ dàng di chuyển người dùng giữa các VLAN mà không cần thay đổi cấu trúc vật lý
    \item \textbf{Hiệu năng:} Giảm congestion bằng cách phân chia traffic
    \item \textbf{Tiết kiệm chi phí:} Không cần mua nhiều switch vật lý cho mỗi nhóm
\end{itemize}

VLAN hoạt động ở Layer 2 (Data Link Layer) của mô hình OSI và được định danh bằng VLAN ID (VID) trong khoảng từ 1 đến 4094. Theo chuẩn IEEE 802.1Q, VLAN ID sử dụng 12 bit, cho phép tối đa 4096 VLAN, trong đó:
\begin{itemize}
    \item VLAN 0: Dành riêng (reserved)
    \item VLAN 1: VLAN mặc định (default VLAN)
    \item VLAN 2-1001: VLAN thông thường (normal range)
    \item VLAN 1002-1005: Dành riêng cho Token Ring và FDDI
    \item VLAN 1006-4094: VLAN mở rộng (extended range)
    \item VLAN 4095: Dành riêng (reserved)
\end{itemize}

\subsubsection{Access Port}
Access Port là loại port trên switch được cấu hình để thuộc về một VLAN duy nhất. Đây là loại port thường được sử dụng để kết nối các thiết bị đầu cuối như máy tính, máy in, điện thoại IP.

\textbf{Đặc điểm của Access Port:}
\begin{itemize}
    \item Chỉ thuộc về một VLAN duy nhất tại một thời điểm
    \item Frame Ethernet không được tag VLAN khi truyền qua access port
    \item Thiết bị đầu cuối không cần biết về sự tồn tại của VLAN
    \item Switch tự động thêm VLAN tag vào frame khi nhận từ access port
    \item Switch tự động gỡ VLAN tag khỏi frame trước khi gửi ra access port
\end{itemize}

\textbf{Ví dụ cấu hình (Cisco):}
\begin{verbatim}
Switch(config)# interface FastEthernet 0/1
Switch(config-if)# switchport mode access
Switch(config-if)# switchport access vlan 10
\end{verbatim}

\subsubsection{Trunk Port}
Trunk Port là loại port được cấu hình để truyền tải traffic của nhiều VLAN cùng lúc. Trunk port thường được sử dụng để kết nối giữa các switch với nhau hoặc kết nối switch với router.

\textbf{Đặc điểm của Trunk Port:}
\begin{itemize}
    \item Có thể truyền traffic của nhiều VLAN
    \item Frame được tag theo chuẩn IEEE 802.1Q (hoặc ISL của Cisco)
    \item Sử dụng khái niệm Native VLAN cho frame không được tag
    \item Bandwidth được chia sẻ giữa các VLAN
\end{itemize}

\textbf{IEEE 802.1Q Frame Format:}

Frame 802.1Q được tạo bởi việc chèn thêm 4 byte VLAN tag vào Ethernet frame chuẩn, ngay sau trường Source MAC Address:

\begin{itemize}
    \item \textbf{TPID (Tag Protocol Identifier) - 2 bytes:} Giá trị 0x8100, định danh frame là 802.1Q tagged frame
    \item \textbf{TCI (Tag Control Information) - 2 bytes:} Gồm 3 trường con:
    \begin{itemize}
        \item \textit{PCP (Priority Code Point) - 3 bits:} Độ ưu tiên của frame (0-7), sử dụng cho QoS
        \item \textit{DEI (Drop Eligible Indicator) - 1 bit:} Cho biết frame có thể bị drop khi congestion
        \item \textit{VID (VLAN Identifier) - 12 bits:} VLAN ID (0-4095)
    \end{itemize}
\end{itemize}

\textbf{Ví dụ cấu hình (Cisco):}
\begin{verbatim}
Switch(config)# interface GigabitEthernet 0/1
Switch(config-if)# switchport mode trunk
Switch(config-if)# switchport trunk allowed vlan 10,20,30
Switch(config-if)# switchport trunk native vlan 99
\end{verbatim}

\textbf{Native VLAN:}
Native VLAN là VLAN mặc định trên trunk port. Frame thuộc Native VLAN không được tag khi truyền qua trunk port. Mặc định, Native VLAN là VLAN 1, nhưng nên thay đổi để tăng cường bảo mật.

\subsubsection{Inter VLAN Routing}
Mặc định, các VLAN khác nhau không thể giao tiếp với nhau vì chúng ở các broadcast domain riêng biệt. Để các VLAN giao tiếp, cần sử dụng Inter-VLAN Routing. Có ba phương pháp chính:

\paragraph{1. Router-on-a-stick (Router với Sub-interface)}
Phương pháp này sử dụng một router vật lý với nhiều sub-interface logic, mỗi sub-interface đại diện cho một VLAN.

\textbf{Nguyên lý hoạt động:}
\begin{itemize}
    \item Router kết nối đến switch qua một trunk link
    \item Tạo sub-interface cho mỗi VLAN trên router
    \item Mỗi sub-interface được gán một IP làm gateway cho VLAN tương ứng
    \item Router nhận frame từ VLAN này, routing, và chuyển đến VLAN khác
\end{itemize}

\textbf{Ưu điểm:}
\begin{itemize}
    \item Chỉ cần một interface vật lý
    \item Tiết kiệm chi phí phần cứng
    \item Dễ cấu hình
\end{itemize}

\textbf{Nhược điểm:}
\begin{itemize}
    \item Băng thông bị giới hạn bởi một interface vật lý
    \item Hiệu năng thấp hơn Layer 3 Switch
    \item Không phù hợp với mạng lớn
\end{itemize}

\textbf{Ví dụ cấu hình:}
\begin{verbatim}
Router(config)# interface GigabitEthernet 0/0
Router(config-if)# no shutdown

Router(config)# interface GigabitEthernet 0/0.10
Router(config-subif)# encapsulation dot1Q 10
Router(config-subif)# ip address 192.168.10.1 255.255.255.0

Router(config)# interface GigabitEthernet 0/0.20
Router(config-subif)# encapsulation dot1Q 20
Router(config-subif)# ip address 192.168.20.1 255.255.255.0
\end{verbatim}

\paragraph{2. Layer 3 Switch (SVI - Switched Virtual Interface)}
Layer 3 Switch là switch có khả năng routing ở Layer 3. Phương pháp này sử dụng SVI (Switched Virtual Interface) - các interface ảo tương ứng với từng VLAN.

\textbf{Nguyên lý hoạt động:}
\begin{itemize}
    \item Tạo interface VLAN (SVI) cho mỗi VLAN
    \item Gán địa chỉ IP cho mỗi SVI làm gateway
    \item Switch thực hiện routing giữa các VLAN bằng ASIC (hardware)
    \item Kích hoạt IP routing trên switch
\end{itemize}

\textbf{Ưu điểm:}
\begin{itemize}
    \item Hiệu năng cao (routing bằng hardware)
    \item Không bị giới hạn băng thông như router-on-a-stick
    \item Phù hợp với mạng lớn
    \item Latency thấp
\end{itemize}

\textbf{Nhược điểm:}
\begin{itemize}
    \item Chi phí cao hơn Layer 2 Switch
    \item Cấu hình phức tạp hơn
\end{itemize}

\textbf{Ví dụ cấu hình:}
\begin{verbatim}
Switch(config)# ip routing

Switch(config)# interface vlan 10
Switch(config-if)# ip address 192.168.10.1 255.255.255.0
Switch(config-if)# no shutdown

Switch(config)# interface vlan 20
Switch(config-if)# ip address 192.168.20.1 255.255.255.0
Switch(config-if)# no shutdown
\end{verbatim}

\paragraph{3. External Router}
Phương pháp này sử dụng một router riêng biệt với nhiều interface vật lý, mỗi interface kết nối đến một VLAN.

\textbf{Nguyên lý hoạt động:}
\begin{itemize}
    \item Mỗi VLAN kết nối đến một interface vật lý riêng trên router
    \item Mỗi interface được gán IP làm gateway cho VLAN
    \item Router thực hiện routing giữa các interface
\end{itemize}

\textbf{Ưu điểm:}
\begin{itemize}
    \item Đơn giản, dễ hiểu
    \item Băng thông riêng biệt cho mỗi VLAN
\end{itemize}

\textbf{Nhược điểm:}
\begin{itemize}
    \item Lãng phí interface
    \item Chi phí cao (nhiều interface, nhiều cable)
    \item Không linh hoạt khi thêm VLAN mới
    \item Ít được sử dụng trong thực tế
\end{itemize}

\textbf{So sánh các phương pháp Inter-VLAN Routing:}

\begin{center}
\begin{tabular}{|l|c|c|c|}
\hline
\textbf{Tiêu chí} & \textbf{Router-on-a-stick} & \textbf{Layer 3 Switch} & \textbf{External Router} \\
\hline
Hiệu năng & Thấp & Cao & Trung bình \\
\hline
Chi phí & Thấp & Cao & Rất cao \\
\hline
Scalability & Thấp & Cao & Thấp \\
\hline
Độ phức tạp & Đơn giản & Trung bình & Đơn giản \\
\hline
Phù hợp & Mạng nhỏ & Mạng lớn/Data Center & Ít sử dụng \\
\hline
\end{tabular}
\end{center}

\subsection{VPN và GRE Tunnel}
\begin{itemize}
    \item \textbf{VPN}(Virtual Private Network) là một mạng riêng ảo cho phép mở rộng kết nối mạng LAN qua một hạ tầng công cộng bằng cách đóng gói và mã hóa các gói dữ liệu.
     \item \textbf{GRE Tunnel} GRE (Generic Routing Encapsulation) là giao thức tạo đường hầm mạng ảo P2P đơn giản.GRE đóng gói hầu hết các loại gói IP (IPv4, IPv6) hoặc các giao thức lớp mạng khác bên trong một gói IP mới. GRE thường dùng để chuyển các lưu lượng đa hướng (multicast, các giao thức định tuyến nhúng) hoặc tạo kênh truyền L2 giữa hai mạng bằng cách gói đa hướng bên trong kênh P2P.
\end{itemize}
